\section{Discussion of Limitations and Future Works}  \label{app:limitations}

Despite the improvements achieved in reaction-aware listening head generation, several limitations remain. 

\textbf{Reaction subjectiveness}. Listener reactions are inherently multi-valid and subjective: given the same conversational cue, different listeners may exhibit different reaction types, timings, or intensities. Our current evaluation follows a single-reference ground-truth setting: the goal is to measure whether the generated listener reaction is consistent with the observed human reaction in the dataset. Therefore, R-F1, R-tIoU, R-ATD, and R-FID should be interpreted as reference-consistency metrics rather than absolute judgments that only one reaction is socially appropriate.

\textbf{Complexity of listener's behaviors}. Our work focuses on a set of explicit \textbf{actual reactions} (e.g., \textit{nodding}, \textit{smiling}), while listener behavior, more broadly, also includes \textbf{micro-behaviors} (e.g., subtle gaze or facial changes) and \textbf{baseline dynamics} (e.g., natural head motion and posture drift showing the person is not stationary). These components are more continuous and harder to annotate, and are not explicitly modeled in our current framework.

\textbf{Multimodal attributes}. Listener behavior is influenced by factors beyond audio prosody, such as language semantics, personality, and conversational context, which are only partially captured in our current model. Extending toward richer multimodal understanding and more comprehensive modeling of listener motion remains an important direction for future work.

\textbf{Preference learning}. We choose supervised temporal reaction modeling instead of reinforcement learning because our goal is to explicitly inject reaction-related temporal supervision into the generation process. The curated dataset provides frame-level reaction annotations, which allow direct optimization of reaction occurrence and temporal consistency. Applying reinforcement learning would require designing a reward function that accurately reflects human judgments of reaction appropriateness, but will theoretically optimize listener's behaviors based on human preferences. Therefore, introducing reinforcement learning will likely further improve listener's generation quality.

\textbf{Full-duplex and multi-party scenarios}. Our GLARE is currently designed to handle listener's motion with reaction generation in conversations with only one speaker, not directly targeting full-duplex or multi-party scenarios, despite being more real-world but complicated at the same time. However, we believe it can be extended to adapt to above tasks. For example, one can adopt any existing full-duplex module to decide when the avatar should speak or listen, in which the corresponding speaking-head module or GLARE would produce facial motion generation based on the switching policy. Furthermore, an interaction-level arbitrator can be introduced to determine the conversational state of each participant in multi-party scenarios, including who is speaking, who is listening, and when state transitions occur. The corresponding speaking or listening generation module can then be activated according to the current interaction state. These would be our future research directions.


\section{Broader Impact}  \label{app:broader_impact}

This work aims to improve the naturalness and appropriateness of listener behavior in conversational agents, with potential applications in virtual assistants, social robots, telepresence, and human-computer interaction systems. By enabling more responsive and context-aware listener reactions, our approach may contribute to more engaging and human-like interactions.

At the same time, improved generation of realistic listener behaviors may introduce risks related to synthetic media, such as creating misleading or deceptive conversational content. In addition, listener reactions can implicitly convey agreement or emotional alignment, which may be misinterpreted or misused in certain contexts. We emphasize that this work is intended for research purposes, and responsible deployment should include appropriate safeguards such as clear disclosure of synthetic content, adherence to dataset licensing terms, and consideration of fairness and diversity in conversational behaviors.